% MT65091C
% Vergiftungen, allgemein
% 14.0
% 2013-11-30
:ref:`m-am-vergiftung`
\begin{enumerate}
  -   Selbstschutz
  -   Beendigung des Giftkontaktes
  -   Evtl. Dekontamination bei
  Gefahrstoffen\footnote{z.\,B. Insektizide, Lösungsmittel
  in der chemischen Industrie, \ldots}
  -   Sicherung der Vitalfunktionen
  -   Entscheidung, ob es sich um einen akut vital
  bedrohten Patienten handelt
  -   Quellenforschung (Mistkübel
  durchsuchen)
  -   Nüchtern lassen
  -   NICHT absichtlich Erbrechen
  herbeiführen
  -   Asservierung (Packungen, Erbrochenes in
  Nierentasse mitnehmen!)
  -   Anamnese, Besonderheiten:
  \begin{multicols}{2}
  \begin{enumerate}\raggedbottom
    -   **Wer**? (Alter, KG, Vorerkrankungen, \ldots
    $\rightarrow$ Risikoabschätzung!)
    -   **Was**? (Art des Giftes $\rightarrow$ evt.
    zu erwartende Symptome)
    -   **Wann**? (Zeit seit Einnahme)
    -   **Wie**?
    (inhalativ, oral, perkutan, intravenös etc.)
    -   **Wieviel**? (Dosisabschätzung)
    -   **Warum**? (Ursachenforschung:
    Selbst\-mord\-ver\-such, Unfall, Fremdverschulden,
    Vorsatz?)
  \end{enumerate}
  \end{multicols}
  -   Ggfs. **Expertenrat einholen** und Entscheidung
  bezüglich des weiteren Vorgehens:
  \begin{multicols}{2}
  \begin{enumerate}\raggedbottom
    -   Gegenmittel verfügbar?
    -   \enquote{Normales} Spitalsbett?
    -   Intensivüberwachung?
    -   Intensivüberwachung mit Entgiftungseinheit?
    -   Dialyse?
    -   Druckkammer?
    -   Sonstige Spezialbehandlungseinheit?
  \end{enumerate}
  \end{multicols}
\end{enumerate}